\documentclass{article}\usepackage{booktabs}\begin{document}
\section*{Tables}

\begin{table}
\centering
\caption{\textbf{Table I.} Data sources. All were accessed on 4 November 2026. AHRF: Area Health Resources Files; PEP: Population Estimates Program; BRFSS: Behavioral Risk Factor Surveillance System; ACS: American Community Survey.}

\begin{tabular}{@{}p{0.29\textwidth}p{0.145\textwidth}p{0.18\textwidth}p{0.31\textwidth}@{}}
\toprule
Source (publisher) & Vintage & Used for & Key variables \\
\midrule
AHRF county file (HRSA)~[16] & 2024--2025 & births; external validation & \texttt{births\_3yr\_avg\_23}; AMA OB/GYN and psychiatrist counts \\
Population estimates (Census)~[17] & Vintage 2024 & universe; women 15--44; births & \texttt{AGE1544\_FEM}; \texttt{BIRTHS2024}; \texttt{POPESTIMATE2024} \\
PLACES county data (CDC)~[14] & 2025 release (BRFSS 2023) & adult depression; mental distress & \texttt{DEPRESSION}, \texttt{MHLTH} crude and age-adjusted \\
Social Vulnerability Index (CDC/ATSDR)~[15] & 2022 (ACS 2018--2022) & social vulnerability & \texttt{RPL\_THEMES}, theme percentiles \\
Rural-Urban Continuum Codes (USDA ERS)~[18] & 2023 & rural classification & RUCC 1--9 \\
Gazetteer, boundaries, ZCTA file (Census)~[19] & 2024; 2020 & frontier density; geocoding & land area; county polygons; ZCTA--county overlap \\
FindTreatment.gov locator API (SAMHSA)~[20] & live directory & treatment listings & type; services incl.\ ``Pregnant/postpartum women'' \\
Health center service delivery sites (HRSA)~[21] & live file & health center capacity & site status, type, location, setting, county \\
NPPES downloadable file V.2 (CMS)~[22] & November 2026 & clinicians & entity type, primary taxonomy, practice ZIP \\
Health Care Provider Taxonomy (NUCC)~[23] & v26.1 & clinician grouping & taxonomy code definitions \\
\bottomrule
\end{tabular}
\end{table}

\begin{table}
\centering
\caption{\textbf{Table II.} Clinician groups by NUCC taxonomy (v26.1) and national counts placed in panel counties.}

\begin{tabular}{@{}p{1.6cm}p{4.6cm}r@{}}
\toprule
Group & Taxonomy (examples of codes) & n \\
\midrule
Obstetric & OB/GYN 207V00000X; Obstetrics 207VX0000X; Maternal-fetal medicine 207VM0101X; Advanced practice midwife 367A00000X; Midwife 176B00000X; NP OB/GYN, women's health, perinatal; CNS women's health, perinatal & 82,517 \\
BH prescriber & Psychiatry 2084P0800X; addiction psychiatry; psychosomatic medicine; addiction medicine (five boards); psychiatric/mental health NP 363LP0808X and CNS & 122,193 \\
BH therapist & Psychologist 103T*; clinical social worker 1041C0700X; counselor 101Y*, mental health, addiction, professional; marriage and family therapist 106H00000X & 1,184,102 \\
Family medicine & Family Medicine 207Q00000X (reported only) & 155,439 \\
\bottomrule
\end{tabular}
\end{table}

\begin{table}
\centering
\caption{\textbf{Table III.} County characteristics by rurality class. BH: behavioral health; PPW SUD: substance use treatment listing with a pregnant/postpartum women program. Scores are national percentile ranks (0--100).}

\begin{tabular}{@{}lrrrr@{}}
\toprule
 & Metro & Rural & Frontier & All \\
\midrule
Counties & 1,164 & 1,563 & 417 & 3,144 \\
Births per year & 3,152,315 & 461,953 & 27,795 & 3,642,063 \\
Share of U.S. births (\%) & 86.6 & 12.7 & 0.8 & 100.0 \\
Women aged 15--44 & 59,278,483 & 7,695,613 & 428,770 & 67,402,866 \\
Adult depression, median (\%) & 23.2 & 23.5 & 20.9 & 23.0 \\
SVI overall, median & 0.47 & 0.56 & 0.34 & 0.50 \\
No PPW SUD program (\% of counties) & 35.1 & 65.5 & 86.6 & 57.0 \\
No behavioral prescriber (\%) & 11.2 & 30.3 & 72.4 & 28.8 \\
No obstetric clinician (\%) & 16.5 & 36.7 & 77.2 & 34.6 \\
No health center site (\%) & 18.6 & 30.6 & 65.9 & 30.8 \\
BH clinicians per 1,000 births, median & 219.5 & 127.8 & 70.0 & 155.4 \\
Obstetric clinicians per 1,000 births, median & 14.3 & 6.8 & 0.0 & 7.9 \\
Need score, median & 49.5 & 56.5 & 35.6 & 51.2 \\
Capacity score, median & 58.8 & 47.7 & 33.8 & 50.4 \\
Gap score, median & −7.9 & 6.9 & −1.6 & −0.1 \\
High need, low capacity (counties) & 23 & 72 & 12 & 107 \\
\bottomrule
\end{tabular}
\end{table}

\begin{table}
\centering
\caption{\textbf{Table IV.} The 15 counties with the largest gap scores. Dep.: adult depression prevalence; PPW: PPW SUD listings; BH: behavioral health clinicians; OB: obstetric clinicians. Need and gap are percentile-based scores.}

\begin{tabular}{@{}rllrrrrrrrr@{}}
\toprule
\# & County & Class & Births & Dep.\,\% & SVI & PPW & BH & OB & Need & Gap \\
\midrule
1 & Mahnomen, MN & rural & 77 & 26.6 & 0.96 & 0 & 2 & 0 & 90.6 & +72.8 \\
2 & Crockett, TN & metro & 164 & 27.9 & 0.72 & 0 & 2 & 0 & 82.2 & +65.3 \\
3 & Lauderdale, TN & rural & 285 & 27.9 & 0.92 & 0 & 9 & 0 & 92.5 & +64.3 \\
4 & Jefferson, OK & rural & 68 & 25.4 & 0.81 & 0 & 0 & 0 & 78.2 & +62.2 \\
5 & McDonald, MO & rural & 326 & 26.5 & 0.89 & 0 & 9 & 0 & 86.8 & +59.8 \\
6 & Newton, TX & rural & 113 & 25.0 & 0.85 & 0 & 4 & 0 & 78.2 & +59.5 \\
7 & Jack, TX & rural & 91 & 24.0 & 0.94 & 0 & 3 & 0 & 77.7 & +59.3 \\
8 & DeKalb, TN & rural & 250 & 30.1 & 0.80 & 0 & 10 & 0 & 88.8 & +59.1 \\
9 & Hopewell city, VA & metro & 321 & 24.1 & 0.88 & 0 & 0 & 0 & 75.1 & +59.1 \\
10 & Rhea, TN & rural & 384 & 30.2 & 0.72 & 0 & 12 & 0 & 85.1 & +58.7 \\
11 & Sequatchie, TN & metro & 181 & 31.3 & 0.50 & 0 & 2 & 0 & 75.0 & +58.2 \\
12 & Henderson, TN & rural & 307 & 29.0 & 0.63 & 0 & 17 & 0 & 78.8 & +58.1 \\
13 & Tyler, TX & rural & 195 & 24.3 & 0.93 & 0 & 11 & 0 & 78.8 & +58.0 \\
14 & St. Mary, LA & rural & 605 & 27.3 & 0.91 & 0 & 20 & 2 & 90.5 & +57.3 \\
15 & Haywood, TN & rural & 193 & 26.7 & 0.89 & 0 & 6 & 0 & 87.4 & +57.1 \\
\bottomrule
\end{tabular}
\end{table}

\begin{table}
\centering
\caption{\textbf{Table V.} Sensitivity of the gap score. ρ: Spearman correlation with baseline gap ranks (among counties scored in both). HNLC n: high-need, low-capacity counties; overlap and Jaccard against the baseline set of 107.}

\begin{tabular}{@{}lrrrr@{}}
\toprule
Specification & ρ & HNLC n & Overlap & Jaccard \\
\midrule
Baseline & 1.000 & 107 & 107 & 1.00 \\
Capacity per 1,000 women aged 15--44 & 0.994 & 104 & 99 & 0.88 \\
Counties with <20 births unranked & >0.999 & 110 & 101 & 0.87 \\
Health centers counted by address & >0.999 & 106 & 104 & 0.95 \\
PPW programs counted by address & >0.999 & 107 & 107 & 1.00 \\
Age-adjusted depression & 0.989 & 115 & 99 & 0.80 \\
Need adds frequent mental distress & 0.975 & 122 & 93 & 0.68 \\
Capacity adds MH listings & 0.973 & 104 & 78 & 0.59 \\
Behavioral prescribers only & 0.973 & 110 & 85 & 0.64 \\
\bottomrule
\end{tabular}
\end{table}

\begin{table}
\centering
\caption{\textbf{Table VI.} Nevada county values. HC: health center sites; other abbreviations as in Table~IV. †Fewer than 20 births.}

\begin{tabular}{@{}llrrrrrrrr@{}}
\toprule
County & Class & Births & Dep.\,\% & SVI & PPW & HC & OB & BH & Gap \\
\midrule
Carson City & metro & 512 & 21.4 & 0.85 & 4 & 1 & 17 & 248 & −20.6 \\
Churchill & frontier & 280 & 22.0 & 0.63 & 1 & 0 & 4 & 104 & −11.0 \\
Clark & metro & 24,629 & 20.3 & 0.86 & 23 & 38 & 414 & 7,126 & −8.8 \\
Douglas & rural & 323 & 18.9 & 0.21 & 0 & 0 & 4 & 92 & −30.0 \\
Elko & frontier & 676 & 21.5 & 0.62 & 1 & 4 & 7 & 77 & −10.9 \\
Esmeralda† & frontier & 2 & 20.2 & 0.91 & 0 & 0 & 0 & 0 & +38.8 \\
Eureka† & frontier & 18 & 23.5 & 0.43 & 0 & 0 & 0 & 0 & +33.2 \\
Humboldt & frontier & 204 & 21.7 & 0.62 & 0 & 0 & 2 & 23 & +14.1 \\
Lander & frontier & 61 & 22.8 & 0.58 & 0 & 1 & 0 & 4 & +12.6 \\
Lincoln & frontier & 52 & 21.2 & 0.57 & 0 & 0 & 0 & 5 & +19.9 \\
Lyon & metro & 622 & 21.4 & 0.56 & 5 & 0 & 1 & 62 & −1.2 \\
Mineral & frontier & 33 & 22.3 & 0.55 & 0 & 0 & 0 & 6 & +19.1 \\
Nye & frontier & 382 & 21.0 & 0.80 & 2 & 3 & 3 & 98 & −17.4 \\
Pershing & frontier & 42 & 19.6 & 0.68 & 0 & 0 & 0 & 6 & +13.5 \\
Storey† & metro & 18 & 19.7 & 0.18 & 0 & 1 & 2 & 21 & −64.8 \\
Washoe & metro & 4,949 & 21.1 & 0.70 & 6 & 12 & 99 & 2,204 & −19.8 \\
White Pine & frontier & 85 & 21.5 & 0.25 & 0 & 0 & 0 & 37 & −7.6 \\
\bottomrule
\end{tabular}
\end{table}
\end{document}
